\documentclass[11pt,a4paper]{article}
\usepackage[T1]{fontenc}
\usepackage{lmodern}
\usepackage[margin=25mm]{geometry}
\usepackage{amsmath,amssymb,amsthm,mathtools}
\usepackage{mathrsfs}
\usepackage{microtype}
\usepackage[dvipsnames]{xcolor}
\usepackage{booktabs}
\usepackage{enumitem}
\usepackage{hyperref}
\hypersetup{colorlinks=true,linkcolor=MidnightBlue,citecolor=MidnightBlue,urlcolor=MidnightBlue,
 pdftitle={Quaternionic Gaussian fibres and exact information loss},
 pdfauthor={Local mathematical contribution; prepared with OpenAI Codex}}
\setlength{\emergencystretch}{2em}
\setlist{nosep}
\numberwithin{equation}{section}
\newtheorem{theorem}{Theorem}[section]
\newtheorem{proposition}[theorem]{Proposition}
\newtheorem{lemma}[theorem]{Lemma}
\theoremstyle{definition}
\newtheorem{definition}[theorem]{Definition}
\theoremstyle{remark}
\newtheorem{remark}[theorem]{Remark}
\newcommand{\HH}{\mathbb H}
\newcommand{\RR}{\mathbb R}
\newcommand{\PP}{\mathcal P}
\newcommand{\Herm}{\operatorname{Herm}}
\newcommand{\Sp}{\operatorname{Sp}}
\newcommand{\Rea}{\operatorname{Re}}
\newcommand{\tr}{\operatorname{tr}}
\newcommand{\diag}{\operatorname{diag}}
\newcommand{\KL}{\operatorname{D}_{\mathrm{KL}}}
\newcommand{\EE}{\mathbb E}
\newcommand{\Cov}{\operatorname{Cov}}
\newcommand{\dd}{\,\mathrm d}
\newcommand{\FI}{\mathcal I}
\newcommand{\Ntwo}{\Delta}
\newcommand{\rtr}{\tau}
\newcommand{\invstar}{^{-\!*}}

\title{Quaternionic Gaussian fibres\\and exact information loss\\[4pt]
\large A constructive measured-bundle example in the original difference coordinates}
\author{Local mathematical contribution\\\small Prepared with OpenAI Codex; not a Nolan-authored manuscript}
\date{4 September 2026 \quad---\quad local review draft}

\begin{document}
\maketitle
\begin{abstract}
For the full fifteen-dimensional cone of positive quaternionic Hermitian
three-by-three matrices, we retain the finite Gaussian measure as well as its
normalized law. Its mass recovers the Moore cubic and its precision-parameter
Fisher metric recovers twice the Hessian of the negative logarithm of that cubic.
For the specific observation $(q_1-q_2,q_2-q_3)$, we prove an explicit
disintegration, including the surface-volume factor on the four-dimensional
affine fibres. The same coordinates yield a closed formula for lost relative
entropy, a complete equality class, and an orthogonal decomposition of all
fifteen Fisher directions into six observed and nine conditional directions.
Every normalization and coordinate change is proved below. These are classical
Gaussian mechanisms worked out as a source-inspired, reproducible example;
no historical novelty or physical field-equation claim is made.
\end{abstract}

\section{Source motivation and precise scope}\label{sec:source}

The measured-fibre viewpoint motivating this note is explicitly credited to
Nolan's paper~\cite{Nolan2025Threads}, with his preferred citation:
\begin{quote}
Nolan, Michael J. (2025). \emph{Threads as Fiber Density, and Holographic
Cosmology}. GitHub.
\end{quote}
Its Definition 2.1 uses a
base-dependent weight multiplying a base measure and a fixed fibre measure;
its weighted projection integrates over that measured total space. The
repository's projection manuscript explicitly proposes Fisher information
and covariance as possible geometric data in its speculative outlook,
Section 9(a)~\cite{NolanRepository2025Projection}. The measured-bundle
framework document presents the broader construction~\cite{NolanRepository2025MBPF}.
We supply a finite-dimensional example in which the weight, conditional law,
Fisher metric and information contraction can all be calculated exactly.
The parameter-dependent trivialization in Section~\ref{sec:disintegration}
exhibits the precise relation to the fixed-fibre product ansatz.

Our Gaussian family is an additional construction, not a claim that those
sources already contain a quaternionic determinant theorem. The source
lineage also includes later anonymously or pseudonymously bylined papers;
we neither use their stochastic or field-theoretic assertions here nor
reassign their bylines to Nolan. The original difference map and marking
used below come from the local quaternionic determining-algebra work,
not from Nolan's PDFs.

Gaussian Fisher geometry and exponential-family Hessian formulas are
standard; see, for example, Nielsen~\cite[Section 2]{Nielsen2023Gaussian}
and Erdmenger, Grosvenor and Jefferson~\cite[Sections 2--2.2]{EGJ2020}.
Our proofs do not assume those formulas. In particular, the metric below is
on the \emph{precision parameters} $H$, not on a location parameter with fixed
covariance, and not on physical spacetime. The new follow-through relative
to the accompanying local bridge is the explicit conditional law and the
sharp KL and Fisher loss under the unchanged observation $P$.

\section{Quaternionic conventions and the original cubic}\label{sec:algebra}

Write $a=a_0+a_1i+a_2j+a_3k\in\HH$, with
$i^2=j^2=k^2=ijk=-1$. Columns form right quaternionic modules; matrices
act on the left. A star denotes conjugate transpose. Lebesgue measure
on $\HH^n$ means the product of its specified $4n$ real coordinate measures.
For a square matrix put $\rtr(A)=\Rea\tr(A)$.
The identity $\Rea(ab)=\Rea(ba)$ follows by expanding real and imaginary
parts. Summing entries proves $\rtr(AB)=\rtr(BA)$, also for rectangular
compatible matrices, hence cyclicity for longer products. Quaternionic
traces without their real part are not used cyclically.

Let $\PP_n=\{H\in\Herm_n(\HH):q^*Hq>0\text{ for }q\ne0\}$.
The quadratic form is real because it equals its conjugate. We retain
the full cone $\PP_3$, with three real diagonal and twelve real
off-diagonal coordinates. For
\begin{equation}\label{eq:cubic}
H=\begin{pmatrix}a&v&r\\\bar v&b&s\\\bar r&\bar s&d\end{pmatrix},
\qquad
N(H)=abd-d|v|^2-b|r|^2-a|s|^2+2\Rea(\bar r v s).
\end{equation}
In rank two use
$\Ntwo\bigl(\begin{smallmatrix}a&v\\\bar v&b\end{smallmatrix}\bigr)=ab-|v|^2$;
in rank one use $N_1(c)=c$. We also write $N_2=\Ntwo$ and $N_3=N$.

The original complex-plus-triplet coordinates are retained by fixing
$\mathbb C=\RR+\RR i\subset\HH$ and setting
\begin{equation}\label{eq:theta}
B(b)=\begin{pmatrix}0&\bar b_3&-\bar b_2\\-\bar b_3&0&\bar b_1\\
\bar b_2&-\bar b_1&0\end{pmatrix},\qquad
\Theta(A,b)=A+B(b)j,
\end{equation}
for $A\in\Herm_3(\mathbb C)$ and $b\in\mathbb C^3$.
The identity
\[
(a+bj)(c+dj)=ac-b\bar d+(ad+b\bar c)j
\]
implies $(A+Bj)^*=A^*-B^Tj$. Thus $\Theta$ is a real-linear
bijection onto $\Herm_3(\HH)$: extract the complex matrix part $A$
and the skew matrix $B$, then recover
$b=(\overline{B_{23}},\overline{B_{31}},\overline{B_{12}})^T$.
In these same coordinates the cubic is exactly
\begin{equation}\label{eq:paircubic}
N(\Theta(A,b))=\det_{\mathbb C}A-b^*Ab.
\end{equation}
For an explicit check, write $A_{23}=z_1$, $A_{31}=z_2$,
$A_{12}=z_3$. The corresponding quaternionic entries are
$h_{23}=z_1+\bar b_1j$, $h_{31}=z_2+\bar b_2j$,
$h_{12}=z_3+\bar b_3j$. The complex part of their product in that order is
\[
z_1z_2z_3-\bar b_1b_2z_3-z_1\bar b_2b_3-\bar b_1\bar z_2b_3.
\]
Its real part agrees with the real triple term in \eqref{eq:cubic}
by real cyclicity. Together with
$|z_i+\bar b_ij|^2=|z_i|^2+|b_i|^2$, its $b$ terms are
\[
-a|b_1|^2-b_{\rm diag}|b_2|^2-d|b_3|^2
-2\Rea(\bar b_1z_3b_2+\bar b_2z_1b_3+\bar b_1\bar z_2b_3),
\]
where $b_{\rm diag}=A_{22}$ distinguishes the diagonal entry from the
triplet $b$. These are exactly $-b^*Ab$; the remaining terms are
$\det A$. This proves \eqref{eq:paircubic} without changing coordinates
or requiring the matrices to commute.

\begin{lemma}[Spectral and real determinant identities]\label{lem:spectral}
Every quaternionic Hermitian matrix is unitarily diagonalizable with real
eigenvalues. On $\PP_n$, $n\le3$, $N_n(H)$ is the positive product of those
eigenvalues. If $\mathscr L(H)$ is its real left-multiplication matrix, then
\begin{equation}\label{eq:realification}
\det_{\RR}\mathscr L(H)=N_n(H)^4,
\qquad \tr_{\RR}\mathscr L(A)=4\rtr(A).
\end{equation}
\end{lemma}
\begin{proof}
The real continuous Rayleigh function $q^*Hq$ has a maximum on the real
unit sphere. At a maximizer $e$, differentiating subject to $e^*e=1$
gives $\Rea(w^*(He-\lambda e))=0$ for every $w\in\HH^n$ and some
real $\lambda$. Taking $w=He-\lambda e$ proves $He=e\lambda$.
The quaternionic orthogonal complement $e^*w=0$ is invariant, since
$e^*Hw=(He)^*w=\lambda e^*w$. Quaternionic Gram--Schmidt followed by
induction gives an orthonormal eigenbasis. Positivity is therefore exactly
positivity of all eigenvalues.

For rank three direct multiplication gives
\begin{align*}
\rtr(H^2)&=a^2+b^2+d^2+2(|v|^2+|r|^2+|s|^2),\\
\rtr(H^3)&=a^3+b^3+d^3+3(a+b)|v|^2+3(a+d)|r|^2\\
&\hspace{16mm}+3(b+d)|s|^2+6\Rea(\bar r v s).
\end{align*}
Substitution into
\[
\frac{\rtr(H)^3-3\rtr(H)\rtr(H^2)+2\rtr(H^3)}6
\]
gives exactly \eqref{eq:cubic}. Cyclicity makes this expression invariant
under unitary conjugation; on a real diagonal it is the product of the
three entries. The analogous identity
$N_2(H)=(\rtr(H)^2-\rtr(H^2))/2$ proves the rank-two assertion.

Explicitly, left multiplication by $a$ has matrix
\[
L(a)=\begin{pmatrix}
a_0&-a_1&-a_2&-a_3\\
a_1&a_0&-a_3&a_2\\
a_2&a_3&a_0&-a_1\\
a_3&-a_2&a_1&a_0
\end{pmatrix}.
\]
Replace each entry of a quaternionic matrix by this block to obtain
$\mathscr L$. Associativity and conjugation give
$\mathscr L(AB)=\mathscr L(A)\mathscr L(B)$ and
$\mathscr L(A^*)=\mathscr L(A)^T$. Its real trace is $4\rtr(A)$.
A unitary quaternionic matrix becomes orthogonal, and a Hermitian
eigenvalue becomes four identical real eigenvalues. This proves
\eqref{eq:realification}.
\end{proof}

For $H>0$ its real positive powers are defined by the proved spectral
decomposition: if $H=U\diag(\lambda_j)U^*$, set
$H^a=U\diag(\lambda_j^a)U^*$ for real $a$.
This is independent of the eigenbasis, since $\lambda^a$ acts as a
real scalar on each eigenspace. In particular $H^{1/2}$ and $H^{-1/2}$
exist and are mutually inverse positive Hermitian matrices.

\begin{lemma}[Differential of the cubic]\label{lem:differential}
For $H\in\PP_n$ and Hermitian tangents $K,L$,
\begin{align}
D\log N_n(H)[K]&=\rtr(H^{-1}K),\label{eq:dlog}\\
D^2(-\log N_n)(H)[K,L]&=\rtr(H^{-1}KH^{-1}L).\label{eq:hessian}
\end{align}
The last expression is a positive-definite symmetric bilinear form.
\end{lemma}
\begin{proof}
For a real matrix $M$, the Leibniz determinant expansion gives
$\det(I+tM)=1+t\tr(M)+O(t^2)$. Factoring an invertible matrix
therefore gives $D\log\det(A)[B]=\tr(A^{-1}B)$ on positive definite
matrices. Apply this to $\log N_n=\tfrac14\log\det_{\RR}\mathscr L(H)$.
Differentiating $H^{-1}H=I$ gives
$D(H^{-1})[L]=-H^{-1}LH^{-1}$, proving \eqref{eq:hessian}.
Real cyclicity gives symmetry. With $T=H^{-1/2}KH^{-1/2}$,
the value at $(K,K)$ is $\rtr(T^2)$, the sum of squared real eigenvalues
of $T$. It is strictly positive for $K\ne0$.
\end{proof}

\section{Finite Gaussian mass, moments and Fisher metric}\label{sec:gaussian}

\begin{definition}
For $H\in\PP_n$, $n\le3$, define the finite measure and probability law
\begin{equation}\label{eq:family}
\mu_H(\mathrm dq)=e^{-q^*Hq}\mathrm d^{4n}q,
\qquad Z_n(H)=\int\mathrm d\mu_H,
\qquad p_H=Z_n(H)^{-1}\mu_H.
\end{equation}
Expectations with subscript $H$ refer to $p_H$. Density and probability
measure will both be denoted $p_H$ when the integration variable is clear.
\end{definition}

\begin{theorem}[Exact mass and precision geometry]\label{thm:gaussian}
With the unchanged coordinate measure of \eqref{eq:family},
\begin{align}
Z_n(H)&=\frac{\pi^{2n}}{N_n(H)^2}, &
\EE_H[qq^*]&=2H^{-1},\label{eq:massmom}\\
\EE_H[q^*Kq]&=2\rtr(H^{-1}K), &
\Cov_H(q^*Kq,q^*Lq)&=2\rtr(H^{-1}KH^{-1}L).\label{eq:quadmom}
\end{align}
The score in direction $K$ and Fisher metric are
\begin{align}
s_{H,K}(q)&=2\rtr(H^{-1}K)-q^*Kq,\label{eq:score}\\
\FI_H(K,L)&=\EE_H[s_{H,K}s_{H,L}]
=2\rtr(H^{-1}KH^{-1}L)
=2D^2(-\log N_n)[K,L].\label{eq:fisher}
\end{align}
In particular, the family is injectively parametrized by $H$.
\end{theorem}
\begin{proof}
The one-dimensional Gaussian integral is $\sqrt\pi$: square it, use
polar coordinates in $\RR^2$, and evaluate
$2\pi\int_0^\infty re^{-r^2}\dd r=\pi$.
Consequently $\int_{\HH}e^{-c|w|^2}\mathrm d^4w=\pi^2/c^2$ for $c>0$.
Diagonalize $H$ by Lemma~\ref{lem:spectral}. Its unitary coordinate
change has absolute real Jacobian one, so multiplying these rank-one
integrals proves the mass formula.

In diagonal coordinates all $4n$ real coordinates are independent,
centered, and each of the four coordinates at eigenvalue $\lambda$
has variance $1/(2\lambda)$. For this last constant, integration by parts
of $(xe^{-\lambda x^2})'$ gives
$\int x^2e^{-\lambda x^2}\dd x=(2\lambda)^{-1}\int e^{-\lambda x^2}\dd x$;
the boundary term vanishes at both infinities. Thus the quaternionic matrix second moment
is diagonal with entries $2/\lambda$. Conjugation back proves
$\EE_H[qq^*]=2H^{-1}$ and real cyclicity proves the first quadratic moment.

To justify all differentiations, choose a neighborhood of $H$ on which
$q^*H'q\ge\epsilon|q|^2$ for some $\epsilon>0$. Each derivative of
the integrand of fixed finite order is bounded by a constant times
a polynomial in $|q|$ times $e^{-\epsilon|q|^2}$, which is integrable
by radial integration. Dominated differentiation is therefore valid.
Differentiating $\log Z_n$ once gives $-\EE_H[q^*Kq]$ and twice gives
the covariance in \eqref{eq:quadmom}. Since
$\log Z_n=2n\log\pi-2\log N_n$, Lemma~\ref{lem:differential}
gives the displayed covariance and Fisher formulas. Differentiation of
the normalized density gives \eqref{eq:score}. Finally,
$H=2(\EE_H[qq^*])^{-1}$ proves injectivity.
\end{proof}

For rank three the finite mass is retained, not discarded by normalization:
\begin{equation}\label{eq:massscaling}
N(H)=\frac{\pi^3}{\sqrt{Z_3(H)}},\qquad
Z_3(rH)=r^{-6}Z_3(H)\quad(r>0).
\end{equation}
The second formula follows either from the degree-three cubic or from
$q\mapsto r^{-1/2}q$, whose twelve-dimensional Jacobian is $r^{-6}$.
The same map pushes $p_H$ to $p_{rH}$.
The determinant-one hypersurface is exactly the mass-$\pi^6$ hypersurface.

\begin{proposition}[Quadratic tilt and its full domain]\label{prop:tilt}
For every Hermitian $K$,
\begin{equation}\label{eq:tilt}
\EE_H e^{q^*Kq}=
\left(\frac{N_n(H)}{N_n(H-K)}\right)^2
\quad\text{if }H-K>0.
\end{equation}
The expectation is infinite otherwise. In the finite case, normalizing
the tilted measure gives exactly $p_{H-K}$.
\end{proposition}
\begin{proof}
Multiplication of densities leaves $e^{-q^*(H-K)q}/Z_n(H)$.
If $H-K>0$, Theorem~\ref{thm:gaussian} gives its mass and normalized law.
Otherwise the spectral lemma gives a zero or negative real eigenvalue.
In its corresponding real coordinate the integrand is constant or grows
as $e^{a x^2}$ with $a>0$. Integration in that coordinate diverges;
Tonelli's theorem applies because the integrand is nonnegative.
\end{proof}

\section{The exact difference observation and disintegration}\label{sec:disintegration}

Keep the original real matrices, acting on quaternionic columns,
\begin{align}
P&=\begin{pmatrix}1&-1&0\\0&1&-1\end{pmatrix},&
u&=\begin{pmatrix}1\\1\\1\end{pmatrix},&
C=PP^*&=\begin{pmatrix}2&-1\\-1&2\end{pmatrix},\label{eq:PU}\\
C^{-1}&=\frac13\begin{pmatrix}2&1\\1&2\end{pmatrix},&
R=P^*C^{-1}&=\frac13\begin{pmatrix}2&1\\-1&1\\-1&-2\end{pmatrix},&
X&=(R\ u).\label{eq:RX}
\end{align}
Direct multiplication gives $PR=I_2$, $Pu=0$, $u^*R=0$, $u^*u=3$,
and hence
\begin{equation}\label{eq:coordinates}
X^{-1}=\begin{pmatrix}P\\u^*/3\end{pmatrix},\quad
z=\binom yt=X^{-1}q,\quad y=Pq,\quad t=u^*q/3,\quad q=Ry+ut.
\end{equation}
Expansion of the real three-by-three determinant gives $\det X=1$.
As a real map on twelve coordinates its determinant is $(\det X)^4=1$;
therefore $\mathrm d^{12}q=\mathrm d^8y\,\mathrm d^4t$.

For $H\in\PP_3$ write
\begin{equation}\label{eq:schur}
B_H=X^*HX=\begin{pmatrix}A_H&b_H\\b_H^*&c_H\end{pmatrix},\quad
\ell_H=c_H^{-1}b_H^*,\quad
S_H=A_H-c_H\ell_H^*\ell_H.
\end{equation}
Here $c_H=u^*Hu>0$ is real, $\ell_H$ is a quaternionic row with two
entries, and $S_H\in\PP_2$. Indeed, keeping the order of quaternionic
products, completion of the square gives
\begin{equation}\label{eq:square}
q^*Hq=y^*S_Hy+c_H|t+\ell_Hy|^2.
\end{equation}
Setting $t=-\ell_Hy$ proves positivity of $S_H$. Conversely, every
$(S,c,\ell)\in\PP_2\times(0,\infty)\times\HH^{1\times2}$ gives
\begin{equation}\label{eq:inverseparam}
B=\begin{pmatrix}S+c\ell^*\ell&c\ell^*\\c\ell&c\end{pmatrix},
\qquad H=X\invstar B X^{-1}\in\PP_3.
\end{equation}
The displayed forward and inverse maps are smooth and inverse to one another.
They are global coordinates of real dimensions $6+1+8=15$; none of
the original parameters has been removed.

\begin{theorem}[Measured fibres, with all volume constants]\label{thm:disintegration}
For the coordinates \eqref{eq:schur},
\begin{align}
N(H)&=c_H\Ntwo(S_H),\label{eq:detfactor}\\
P_*\mu_H(\mathrm dy)&=\lambda_H(y)\mathrm d^8y,
&\lambda_H(y)&=\frac{\pi^2}{c_H^2}e^{-y^*S_Hy},\label{eq:weightedmarginal}\\
P_*p_H&=p_{S_H},&
p_{S_H}(y)&=\frac{\Ntwo(S_H)^2}{\pi^4}e^{-y^*S_Hy}.
\label{eq:marginal}
\end{align}
The conditional probability in the $t$ coordinate is
\begin{equation}\label{eq:conditional}
\kappa_H(y,\mathrm dt)=k_H(t\mid y)\mathrm d^4t,
\qquad k_H(t\mid y)=\frac{c_H^2}{\pi^2}e^{-c_H|t+\ell_Hy|^2}.
\end{equation}
Thus for every nonnegative Borel $f$, and also every absolutely
$\mu_H$-integrable $f$,
\begin{equation}\label{eq:disintegration}
\int f(q)\dd\mu_H(q)
=\int_{\HH^2}\lambda_H(y)
 \left[\int_{\HH}f(Ry+ut)\,\kappa_H(y,\mathrm dt)\right]\mathrm d^8y.
\end{equation}
On the actual fibre $P^{-1}(y)$ the Euclidean surface measure is
$\mathrm d\mathcal H_y^4=9\mathrm d^4t$. Relative to it the conditional
density is $k_H(t\mid y)/9$.
\end{theorem}
\begin{proof}
The rank-one Gaussian integral in $t+\ell_Hy$ equals $\pi^2/c_H^2$,
which proves \eqref{eq:weightedmarginal}. Integrating in $y$ gives
$Z_3(H)=\pi^6/(c_H^2\Ntwo(S_H)^2)$. Comparison with
\eqref{eq:massmom}, and positivity of all factors, proves
\eqref{eq:detfactor}. This also proves \eqref{eq:marginal} after
normalizing by $Z_3(H)$. The density \eqref{eq:conditional} has mass one
by translation and the same rank-one integral. Multiplication of
\eqref{eq:weightedmarginal} and \eqref{eq:conditional} returns exactly
the exponential in \eqref{eq:square}. The coordinate Jacobian and
Tonelli/Fubini therefore prove \eqref{eq:disintegration} with no implicit
disintegration assumption.

The tangent map of $t\mapsto Ry+ut$ sends each real coordinate vector
to three identical copies. Its real Gram matrix is $3I_4$, so the
four-dimensional surface Jacobian is $\sqrt{\det(3I_4)}=9$.
Equivalently the real observation matrix $P_{\RR}$ satisfies
\[
\sqrt{\det_{\RR}(P_{\RR}P_{\RR}^T)}
=\sqrt{\det(C\otimes I_4)}=\det(C)^2=9.
\]
Thus $\mathrm d^{12}q=\tfrac19\mathrm d^8y\,\mathrm d\mathcal H_y^4$
and the asserted conditional density follows.
\end{proof}

\begin{proposition}[Conditional moments and the fixed-fibre presentation]\label{prop:conditionalmom}
Let $m_H(y)=(R-u\ell_H)y$. Then
\begin{align}
S_H^{-1}&=PH^{-1}P^*, &
\EE_H[yy^*]&=2S_H^{-1},\label{eq:marginalcov}\\
m_H(y)&=H^{-1}P^*S_Hy, &
\EE_H[q\mid y]&=m_H(y),\label{eq:condmean}\\
\EE_H[t\mid y]&=-\ell_Hy,&
\EE_H[|t+\ell_Hy|^2\mid y]&=2/c_H,\label{eq:tmom}\\
\EE_H[(q-m_H(y))(q-m_H(y))^*\mid y]&=\frac2{c_H}uu^*.
\label{eq:condcov}
\end{align}
The coordinate $\eta=\sqrt{c_H}(t+\ell_Hy)$ is independent of $y$
under $p_H$, with fixed density $\pi^{-2}e^{-|\eta|^2}\mathrm d^4\eta$.
In these coordinates the finite measure is exactly
\begin{equation}\label{eq:fixedfibre}
\mathrm d\mu_H
=\lambda_H(y)\mathrm d^8y\,
 \bigl(\pi^{-2}e^{-|\eta|^2}\mathrm d^4\eta\bigr).
\end{equation}
\end{proposition}
\begin{proof}
The factorization
\[
B_H=\begin{pmatrix}I&\ell_H^*\\0&1\end{pmatrix}
\begin{pmatrix}S_H&0\\0&c_H\end{pmatrix}
\begin{pmatrix}I&0\\\ell_H&1\end{pmatrix}
\]
gives the full inverse
\[
B_H^{-1}=\begin{pmatrix}
S_H^{-1}&-S_H^{-1}\ell_H^*\\
-\ell_HS_H^{-1}&c_H^{-1}+\ell_HS_H^{-1}\ell_H^*
\end{pmatrix}.
\]
Since $H^{-1}=XB_H^{-1}X^*$ and $PX=(I\ 0)$, its upper projected
block is $S_H^{-1}$. Also
$H^{-1}P^*S_H=X\binom I{-\ell_H}=R-u\ell_H$.
The marginal covariance follows from Theorem~\ref{thm:gaussian} in
rank two. The centered conditional quaternion $w=t+\ell_Hy$ has
four independent real coordinates of variance $1/(2c_H)$; it has mean
zero and $\EE|w|^2=2/c_H$. This proves all conditional moments,
including \eqref{eq:condcov} because $q-m_H(y)=uw$.
Finally $\mathrm d^4\eta=c_H^2\mathrm d^4t$. Substitution in
the joint density gives \eqref{eq:fixedfibre}, and after normalization
it is the product of $p_{S_H}$ and the fixed $\eta$ density.
\end{proof}

Formula \eqref{eq:fixedfibre} is the promised exact morphism to the
weighted fixed-fibre form motivating the construction
\cite{Nolan2025Threads,NolanRepository2025MBPF}. The base density is
$\lambda_H$, while the normalization and conditional shape have not been
confused. On an affine fibre the surface-density measure
$e^{-q^*Hq}\mathrm d\mathcal H_y^4$ has mass $9\lambda_H(y)$.
It is \emph{not} the literal restriction of the twelve-dimensional
measure $\mu_H$ to that four-plane: that restriction is zero, since
the plane has twelve-dimensional Lebesgue measure zero. The conditional
law can equivalently be written on the plane as
\begin{equation}\label{eq:surfacedensity}
\frac{c_H^2}{9\pi^2}
\exp\!\left(-\frac{c_H}{3}|q-m_H(y)|^2\right)
\mathrm d\mathcal H_y^4(q),
\end{equation}
because $|uw|^2=3|w|^2$.

\section{Relative entropy and the sharp equality class}\label{sec:KL}

For positive densities set $\KL(p\|r)=\int p\log(p/r)$.
All log-density ratios below are a constant plus a quadratic polynomial,
so their absolute integrability follows from Theorem~\ref{thm:gaussian}.

\begin{proposition}[Quaternionic Gaussian relative entropy]\label{prop:KL}
For $H,G\in\PP_n$, $n\le3$,
\begin{equation}\label{eq:KL}
\KL(p_H\|p_G)
=2\left[\rtr(GH^{-1})-n-\log\frac{N_n(G)}{N_n(H)}\right]
=2\sum_{j=1}^n(\lambda_j-1-\log\lambda_j),
\end{equation}
where $\lambda_j>0$ are the eigenvalues of $H^{-1/2}GH^{-1/2}$.
It is nonnegative, with equality exactly when $G=H$.
\end{proposition}
\begin{proof}
The log-density ratio is
$2\log(N_n(H)/N_n(G))+q^*(G-H)q$.
Taking expectation gives the first equality. Real cyclicity gives
$\rtr(GH^{-1})=\sum_j\lambda_j$.
For the determinant term, multiplicativity of real determinants and
\eqref{eq:realification} give
\[
N_n(H^{-1/2}GH^{-1/2})^4=N_n(H)^{-4}N_n(G)^4.
\]
Positive fourth roots and Lemma~\ref{lem:spectral} give
$\prod_j\lambda_j=N_n(G)/N_n(H)$, proving the second equality.
The scalar function $x-1-\log x$ decreases to zero on $(0,1]$ and
increases from zero on $[1,\infty)$, since its derivative is $1-1/x$.
Hence equality forces every $\lambda_j=1$; the Hermitian spectral
decomposition then gives $H^{-1/2}GH^{-1/2}=I$ and $G=H$.
\end{proof}

\begin{theorem}[Exact KL loss and all equality cases]\label{thm:KLloss}
For the fixed observation $P$ of \eqref{eq:PU}, set
$\delta\ell=\ell_G-\ell_H$. Then
\begin{align}
\KL(p_H\|p_G)
={}&\KL(p_{S_H}\|p_{S_G})
 +2\left[\frac{c_G}{c_H}-1-\log\frac{c_G}{c_H}\right]
 +2c_G\delta\ell S_H^{-1}\delta\ell^*.
\label{eq:KLloss}
\end{align}
Thus relative entropy cannot increase under $P$. Equality holds if and
only if each of the following equivalent conditions holds:
\begin{align}
&c_G=c_H\quad\text{and}\quad\ell_G=\ell_H;\label{eq:equalitycl}\\
&(G-H)u=0;\label{eq:equalityu}\\
&G-H=P^*TP\quad\text{for a unique }T\in\Herm_2(\HH).
\label{eq:equalityT}
\end{align}
In this case $T=S_G-S_H=R^*(G-H)R$. The exact equality class through
$H$ is
\begin{equation}\label{eq:affineslice}
\bigl(H+P^*\Herm_2(\HH)P\bigr)\cap\PP_3.
\end{equation}
\end{theorem}
\begin{proof}
The joint probability density is $p_{S_H}(y)k_H(t\mid y)$.
Split the logarithm of its ratio to the corresponding $G$ density
into marginal and conditional logarithms and use Fubini. Under the
$H$ conditional law, $w=t+\ell_Hy$ has mean zero and
$\EE[|w|^2\mid y]=2/c_H$. Therefore
\begin{align*}
\int k_H(t\mid y)\log\frac{k_H(t\mid y)}{k_G(t\mid y)}\mathrm d^4t
&=2\log\frac{c_H}{c_G}
 +\EE[-c_H|w|^2+c_G|w+\delta\ell y|^2\mid y]\\
&=2\left[\frac{c_G}{c_H}-1-\log\frac{c_G}{c_H}\right]
 +c_G|\delta\ell y|^2.
\end{align*}
The cross term vanishes by the centered conditional mean. Moreover,
\[
\EE_H|\delta\ell y|^2
=\delta\ell\,\EE_H[yy^*]\,\delta\ell^*
=2\delta\ell S_H^{-1}\delta\ell^*.
\]
This proves \eqref{eq:KLloss}. The scalar term is nonnegative by the
same one-variable inequality as in Proposition~\ref{prop:KL}; the
last term is nonnegative because $S_H^{-1}>0$. Their simultaneous
vanishing is exactly \eqref{eq:equalitycl}.

For the linear characterization put $E=G-H$. The last column of
$X^*EX$ is $X^*Eu$; since $X$ is invertible it vanishes exactly when
$Eu=0$. Hermitian symmetry then also kills the last row. Thus
$Eu=0$ is equivalent to
$X^*EX=\diag(T,0)$ for some $T\in\Herm_2(\HH)$.
Transforming back by \eqref{eq:coordinates} gives $E=P^*TP$.
Multiplication by $R^*$ and $R$ proves uniqueness and
$T=R^*ER$. The same zero last row and column say $b_G=b_H$ and
$c_G=c_H$, which is equivalent to \eqref{eq:equalitycl}; in that case
the top Schur blocks differ by $T=S_G-S_H$. This proves every
equivalence and \eqref{eq:affineslice}.
\end{proof}

The marginal term in \eqref{eq:KLloss} is explicitly
\[
2\left[\rtr(S_GS_H^{-1})-2-
\log\frac{\Ntwo(S_G)}{\Ntwo(S_H)}\right].
\]
Within the slice \eqref{eq:affineslice}, the conditional kernel
\eqref{eq:conditional} is independent of the varying parameter $T$.
Consequently $Pq$ is sufficient for that family in the concrete sense
that the same conditional kernel reconstructs every full law from its
marginal. Conversely, any family containing $H$ whose pairwise KL loss
under $P$ vanishes lies in this slice, by Theorem~\ref{thm:KLloss}.
This is a maximal exact equality family, not a sufficiency assertion
for the whole fifteen-dimensional model.

\section{The full Fisher split and infinitesimal equality}\label{sec:fishersplit}

For a tangent $K\in\Herm_3(\HH)$ write dots for derivatives of
\eqref{eq:schur} along $H+\varepsilon K$ at $\varepsilon=0$.
If $X^*KX=\bigl(\begin{smallmatrix}\dot A&\dot b\\\dot b^*&\dot c\end{smallmatrix}\bigr)$,
the exact coordinate derivatives are
\begin{equation}\label{eq:tangents}
\dot\ell=\frac{\dot b^*-\dot c\ell}{c},\qquad
\dot S=\dot A-\dot c\ell^*\ell
 -c\dot\ell^*\ell-c\ell^*\dot\ell.
\end{equation}
Unadorned $S,c,\ell$ in this section denote the coordinates of $H$.

\begin{theorem}[Six observed and nine conditional directions]\label{thm:Fisherloss}
For tangents $K,L$ with coordinate derivatives
$(\dot S_K,\dot c_K,\dot\ell_K)$ and
$(\dot S_L,\dot c_L,\dot\ell_L)$,
\begin{align}
\FI_H(K,L)={}&2\rtr(S^{-1}\dot S_K S^{-1}\dot S_L)
 +\frac{2\dot c_K\dot c_L}{c^2}
 +4c\Rea(\dot\ell_K S^{-1}\dot\ell_L^*).
\label{eq:fishersplit}
\end{align}
The first term is the Fisher metric of the observed family $p_S$.
The remaining terms are the exact information lost by observing $Pq$.
In particular,
\begin{equation}\label{eq:fisherloss}
\FI_H(K,K)-\FI_S(\dot S_K,\dot S_K)
=2(\dot c_K/c)^2+4c\dot\ell_K S^{-1}\dot\ell_K^*\ge0.
\end{equation}
Equality holds exactly when
\begin{equation}\label{eq:fisherzero}
\dot c_K=0,\quad\dot\ell_K=0
\quad\Longleftrightarrow\quad Ku=0
\quad\Longleftrightarrow\quad K=P^*\dot S_KP.
\end{equation}
\end{theorem}
\begin{proof}
At the fixed base point set
$Q=\bigl(\begin{smallmatrix}I&0\\-\ell&1\end{smallmatrix}\bigr)$.
Direct block multiplication using \eqref{eq:tangents} gives
\begin{equation}\label{eq:whitenedblocks}
Q^*B_HQ=\begin{pmatrix}S&0\\0&c\end{pmatrix}=F,
\qquad
Q^*(X^*KX)Q
=\begin{pmatrix}\dot S_K&c\dot\ell_K^*\\c\dot\ell_K&\dot c_K\end{pmatrix}=M_K.
\end{equation}
Here $Q$ is evaluated at $H$; we are transforming tangents at that
point, not differentiating $Q$ a second time. Congruence by an invertible
matrix preserves $\rtr(H^{-1}KH^{-1}L)$, as is verified by substituting
the inverse of the congruence and cancelling adjacent factors under
the real cyclic trace. Thus \eqref{eq:fisher} equals
$2\rtr(F^{-1}M_KF^{-1}M_L)$.
The two diagonal blocks in this product contribute
\[
\rtr(S^{-1}\dot S_KS^{-1}\dot S_L)
+\dot c_K\dot c_L/c^2
+c\Rea(\dot\ell_LS^{-1}\dot\ell_K^*)
+c\Rea(\dot\ell_KS^{-1}\dot\ell_L^*).
\]
The last two real scalars agree by conjugation. This proves
\eqref{eq:fishersplit}. The rank-two case of Theorem~\ref{thm:gaussian}
identifies the marginal term. Positivity of $S^{-1}$ proves the
inequality and forces $\dot c=\dot\ell=0$ at equality.
By \eqref{eq:tangents} these two conditions are exactly $\dot c=\dot b=0$,
or vanishing of the last column of $X^*KX$. The same argument as in
Theorem~\ref{thm:KLloss} gives $Ku=0$ and
$K=P^*\dot S P$, completing all equivalences.
\end{proof}

Because \eqref{eq:inverseparam} is a global coordinate chart, the
$S$ directions have real dimension six and the $(c,\ell)$ directions
have dimensions one and eight. Formula \eqref{eq:fishersplit} says these
three coordinate subspaces are mutually orthogonal at every point.
It does not assert that the metric is a product of three constant
metrics: the $\ell$ block depends on both $S$ and $c$.
The link to finite loss is also exact to second order:
\begin{equation}\label{eq:localKL}
\KL(p_H\|p_{H+\varepsilon K})
=\tfrac12\varepsilon^2\FI_H(K,K)+O(\varepsilon^3).
\end{equation}
Indeed \eqref{eq:dlog} cancels the linear term in \eqref{eq:KL};
\eqref{eq:hessian} supplies the quadratic term, while smoothness of
$\log N$ on the positive cone gives the Taylor remainder. Expanding
\eqref{eq:KLloss} yields the same residual quadratic form as
\eqref{eq:fisherloss}, with the factor $1/2$ in \eqref{eq:localKL}.

\section{Marking, vacuum and gauge transport}\label{sec:gauge}

The observation coordinates $z=X^{-1}q$ must not be confused with the
original latent marking. Retain the latter exactly as
\begin{equation}\label{eq:marking}
D=X^*X=\diag(C^{-1},3),\quad
\Psi=X^{-1}HX\invstar,\quad
\xi=X^*q=Dz.
\end{equation}
Consequently $H=X\Psi X^*$, $q^*Hq=\xi^*\Psi\xi$, and
\begin{equation}\label{eq:Bpsi}
B_H=D\Psi D.
\end{equation}
These identities follow by direct substitution and preserve both the
quadratic form and the coordinate measure: $\det X=1$ also gives
$\mathrm d^{12}\xi=\mathrm d^{12}q$.
It follows from Theorem~\ref{thm:gaussian} that
$\EE_H[\xi\xi^*]=2\Psi^{-1}$.
At the unchanged scalar vacuum $H=vI_3$, $v>0$, the entire dictionary is
\begin{equation}\label{eq:vacuum}
\begin{gathered}
\Psi=vD^{-1},\qquad B_H=vD,\qquad
S_H=vC^{-1},\quad c_H=3v,\quad\ell_H=0,\\
\EE_H[yy^*]=2C/v,\qquad
\EE_H|t|^2=2/(3v),\qquad
\EE_H[\xi\xi^*]=2D/v.
\end{gathered}
\end{equation}

\begin{proposition}[Fixed marking and co-transported observation]\label{prop:gauge}
Let $U\in\Sp(3)=\{U:U^*U=I\}$, $q_U=Uq$ and $H_U=UHU^*$.
Then $\mu_{H_U}=U_*\mu_H$ and $p_{H_U}=U_*p_H$.
With $X$ fixed, put $U'=X^{-1}UX$. The marked transformations are
\begin{equation}\label{eq:fixedgauge}
\Psi_U=U'\Psi U'^*,\qquad
\xi_U=(U'^*)^{-1}\xi,\qquad U'^*DU'=D.
\end{equation}
If instead the observation and its splitting are co-transported by
\begin{equation}\label{eq:transport}
P_U=PU^*,\qquad R_U=UR,\qquad u_U=Uu,\qquad X_U=UX,
\end{equation}
then $y_U=y$, $t_U=t$, $z_U=z$, $D_U=D$, $\xi_U=\xi$,
$\Psi_U=\Psi$, and $B_{H_U}=B_H$, where every marked expression now
uses $X_U$. All the conditional and information identities are
unchanged under this simultaneous transport.
\end{proposition}
\begin{proof}
Unitarity gives $q_U^*H_Uq_U=q^*Hq$ and a real orthogonal Jacobian of
absolute value one. Thus it pushes both measures as asserted.
For fixed $X$, insert $H_U$ into $X^{-1}H_UX\invstar$ to obtain the
first formula of \eqref{eq:fixedgauge}. Also
\[
X^*Uq=X^*UX\invstar\xi=(U'^*)^{-1}\xi,
\]
where $U^{-1}=U^*$. Substitution of $D=X^*X$ proves
$U'^*DU'=D$. In particular, the marked quadratic form is invariant
under the congruence/contragredient pair in \eqref{eq:fixedgauge}.

For co-transport, $P_Uq_U=Pq$, $u_U^*q_U/3=u^*q/3$ and
$X_U^{-1}q_U=X^{-1}q$. Likewise
$X_U^*q_U=X^*q$, $X_U^*X_U=D$,
$X_U^{-1}H_UX_U\invstar=\Psi$, and
$X_U^*H_UX_U=B_H$. Each coordinate and law is therefore identical.
The real absolute Jacobian of $X_U$ remains one since that of $U$ is one.
\end{proof}

For clarity, a \emph{fixed} $P$ is not invariant under all of $\Sp(3)$.
Take $H=\diag(1,2,3)$ and let $U$ interchange the first two coordinate
vectors. Direct inversion and multiplication give
\begin{equation}\label{eq:fixedcounterexample}
PH^{-1}P^*=\begin{pmatrix}3/2&-1/2\\-1/2&5/6\end{pmatrix},\qquad
PH_U^{-1}P^*=\begin{pmatrix}3/2&-1\\-1&4/3\end{pmatrix}.
\end{equation}
The projected covariances, twice these matrices, differ. The
co-transport \eqref{eq:transport}, not an unproved invariance of the
fixed observation, is the exact equivariant statement.

\section{Reproducibility and contribution boundary}\label{sec:verification}

The accompanying exact-arithmetic program
\texttt{verification/check\_gaussian\_projection.py} implements quaternionic
left multiplication independently in real blocks. It checks the displayed
coordinate maps, cubic/determinant relation, Schur inverse and conditional
covariance identities, all fifteen real Hermitian tangent coordinates
and all 225 bilinear Fisher pairings, the algebraic parts of the KL
chain rule, its equality subspace, and the noninvariant fixed-$P$ example.
The report identifies the exact rational test points used. Finite tests
are executable error checks, not proofs of the universally quantified
theorems; the proofs are the preceding text. Build and execution
instructions are supplied with the source.

The proposed contribution is a mathematical example for the measured-bundle
framework, with Nolan's density paper cited in both the source motivation
and the fixed-fibre construction. It neither establishes source claims
about stationarity or helicity nor modifies the separate local scalar or
Yang--Mills actions. No quantum mass gap, physical dynamics, or historical
priority is inferred from the Fisher metric. A prospective upstream
submission should retain the proofs, verification, original coordinate
dictionary and attribution together.

\input{bibliography.tex}
\end{document}
